\documentclass{article}
\usepackage{graphicx}
\usepackage{amsfonts}
\usepackage{amsmath}
\usepackage{times}
\usepackage{lipsum}
\usepackage{xcolor}
\usepackage{tikz}

\title{Meu primeiro documento}
\author{eduardo.salles}
\date{October 2023}

\begin{document}

\pagecolor{yellow! 10}
\newcommand{\institution}{IFIRN}
\renewcommand{\epsilon}{\varepsilon}
\maketitle

\section{Introduction}

\textcolor{red}{\lipsum[1]}

\bigbreak

\colorbox{gray}{grifando}

\newpage

\section{usando o pikz}

\begin{tikzpicture}
    \draw (0,0) -- (1,1);
\end{tikzpicture}

\newpage

\section{macros}

A instituição de ensino é {\institution}

\subsection{alterando comandos já existentes}

Usando o novo comando para o $\epsilon > 0$

\end{document}